\section{Bounded Replay QA and Artifact-Integrity Examples}
\label{sec:oc133-bounded-replay-qa-artifact-integrity-examples}
This section states the empirical role of the 1.3.3 release in prose. The release contains bounded replay QA and artifact-integrity examples, not a completed empirical conquest of every domain. The distinction matters: a target-blind replay row demonstrates that a claim can be tied to a formula, a source snapshot, a comparator, a residual calculation, a negative control, a falsifier, and a replay hash. It does not by itself license a broader domain theory.

Physics rows are read as constants-and-reconstruction checks. A reader should inspect the pinned source, the formula, the observed value, the uncertainty or tolerance, the comparator baseline, and the falsifier. The scientific question is whether the release can preserve the difference between a replayable bounded statement and a grand physics claim. If the row loses that difference, promotion is too strong.

Chemistry rows are read in the same way. A PubChem or NIST-derived row is not a claim that OC has replaced chemistry. It is a claim that the release machinery can bind a chemical statement to a source, formula, comparator, residual, negative control, and falsifier without drifting into unsupported chemical novelty.

Biology rows are more sensitive because biological interpretation is easy to overread. The bounded role is therefore explicit: the row demonstrates replay discipline and model-card routing for a named biological signal or count. It does not claim that the release closes organismal biology, evolution, development, or systems biology as completed empirical domains.

Systems and civilizational rows are bounded even more tightly. A World Bank or macro-series replay can show snapshot integrity, target-blind reconstruction, comparator behavior, residual accounting, and reopening conditions. It cannot, in this release, become a scoped social-science forecast or a completed model of civilizational dynamics.

Mathematics rows are not empirical rows in the same sense. They anchor executable finite-model behavior and proof-route discipline. The important question is whether the finite case was computed independently from raw model facts and whether the tamper controls fail when they should. A label-only success is not evidence.

The comparator baseline is part of the evidence, not an optional column. Without a comparator, a residual is hard to interpret; without a residual, a prediction is not measured; without a negative control, success may be self-confirmation; without a falsifier, the claim cannot be scientifically reopened. The release requires all of these pieces before a replay row can support even bounded public wording.

The replay hash gives the row its reproducibility perimeter. A reader or auditor should be able to ask whether the same source snapshot and same reconstruction script produce the same table. If not, the row is not silently massaged into compatibility; the artifact is regenerated, the mismatch is explained, or the public claim is demoted.

This is why the release separates model-core promotion from empirical ambition. The typed model can be defended by formal and finite evidence while empirical domains continue to mature. That separation is not weakness; it is what prevents the release from pretending that replay QA is already full domain closure.

A hostile reviewer should therefore attack each empirical row by asking five questions: is the source pinned, is the formula explicit, is the comparator meaningful, is the residual accounted for, and is the falsifier capable of failing the row? A row that cannot answer those questions is not eligible for promoted public wording in this release.

An editor should ask a different question: can a reader understand the role of the row without reading the raw JSON? The monograph must answer yes. The raw table remains in the evidence package for reproducibility, while the prose explains what the row does and what it does not do.

Future releases may expand these lanes into stronger empirical claims. The 1.3.3 standard is that expansion must happen through new source locks, held-out or target-blind protocols, comparators, uncertainty, residuals, negative controls, falsifiers, replay hashes, and adversarial review. It cannot happen through rhetorical confidence alone.


\section{Promoted Model-Core Claims in Prose}
The claim register is translated into prose so that a reader can see what is promoted, what is bounded, and what remains research. The register is not allowed to promote unsupported total finality.

This section turns selected evidence records into a readable scientific route. The main prose states what is being claimed, what kind of support carries it, and what condition would reopen it. Complete records remain in the evidence package so that the monograph can stay argumentative rather than archival.

\subsection{Lifecycle status: liveness, death, residue, and rebirth}
The scientific reading is as follows. Evidence support: Death blocks live status; residue and rebirth are token-bound evidence relations with distinct class-specific endpoint rules: residue separates source from residue while returning to the source endpoint, and rebirth separates source, residue, and new target tokens. Rebirth is non-identity unless endpoint-bound identity evidence has identity class, declared i. The evidence package records the complete field. Lean subset, structured proof sheet, and finite witness public proof sheet for the lifecycle status theorem Bounded to stated theorem assumptions, finite witnesses, and public falsifier boundary. The release consequence is local: this entry supports only the stated claim under its declared assumptions and can be reopened at the named boundary.

\subsection{Continuumness-zero obstruction}
The evidence interpretation is deliberately bounded. Evidence support: Lean subset, structured proof sheet, and finite witness public proof sheet for the K-zero boundary theorem Bounded to stated theorem assumptions, finite witnesses, and public falsifier boundary. Recorded content: Continuumness zero requires live support, an independently clear obstruction register, and a declared zero-cause family; a zero-cause label alone does not compute k=0. The important distinction is between the public statement, the artifact that carries it, and the condition that would make the statement too strong.

\subsection{Metric-threshold boundary specialization}
The reviewer route starts from support, not from assertion. Evidence support: Lean subset, structured proof sheet, and finite witness public proof sheet for the boundary representation theorem Bounded to stated theorem assumptions, finite witnesses, and public falsifier boundary. Recorded content: Metric thresholds are a specialization of typed classifier boundaries. A hostile reader should test the support class first, then ask whether the wording stays within that support class.

\subsection{Typed hybrid update and chart-labelled operator semantics}
The boundary reading prevents overclaim. Evidence support: OC operators are typed update semantics; chart-labelled flow-one notation is admitted only for declared chart records, while proof/rewrite and guard/reset updates remain first-class non-smooth cases. No differentiability or ODE-solution theorem is promoted in current formal profile. Lean subset, structured proof sheet, and finite witness public proof sheet for the hybrid semantics theorem Bounded to stated theorem assumptions, finite witnesses, and public falsifier boundary. This entry is included because it constrains promotion; it does not license a stronger conclusion than the proof, finite case, replay row, or comparator can carry.

Synthesis checkpoint 4. The preceding entries should be read together: promotion is earned by an alignment between claim wording, evidence class, and reopening rule. If any one of those three moves, the public sentence has to move with it.

\subsection{Live-status cycle-mode requirement}
The scientific reading is as follows. Evidence support: Lean subset, structured proof sheet, and finite witness public proof sheet for the cycle-mode theorem Bounded to stated theorem assumptions, finite witnesses, and public falsifier boundary. Recorded content: Declared eligible-live status requires an explicit cycle mode or non-vacuous maintenance predicate. The release consequence is local: this entry supports only the stated claim under its declared assumptions and can be reopened at the named boundary.

\subsection{Identity, residue, and rebirth classification}
The evidence interpretation is deliberately bounded. Evidence support: Identity continuation requires endpoint-bound identity evidence: identity class, declared invariant preservation, no residue token, equal source/target endpoint evidence, lifecycle identity-invariant truth, and typed source/target binding. Residue and rebirth evidence classes do not become identity continuation merely by preserving some invariants. No categorical Hom/composition theorem is promoted. Lean subset, structured proof sheet, and finite witness public proof sheet for the identity and rebirth theorem Bounded to stated theorem assumptions, finite witnesses, and public falsifier boundary. The important distinction is between the public statement, the artifact that carries it, and the condition that would make the statement too strong.

\subsection{Declared component independence of the semantic verdict interface}
The reviewer route starts from support, not from assertion. Evidence support: Within the declared current release tuple semantics, each tuple component has a one-field semantic keep/drop witness that changes the release verdict. Lean subset, structured proof sheet, and finite witness public proof sheet for the minimality witness theorem Bounded to stated theorem assumptions, finite witnesses, and public falsifier boundary. A hostile reader should test the support class first, then ask whether the wording stays within that support class.

\subsection{Adjacent K-level witness consistency}
The boundary reading prevents overclaim. Evidence support: Every declared adjacent K-level transition K0-\textgreater{}K12 has a release-atlas row, retained-witness evaluator check, executable finite row, and inert-witness demotion control inside the current formal profile release classifier; independent semantic irreducibility beyond this declared classifier is a future proof obligation, not a promoted current formal profile theorem. Lean subset, structured proof sheet, and finite witness public proof sheet for the K-level witness theorem Bounded to stated theorem assumptions, finite witnesses, and public falsifier boundary. This entry is included because it constrains promotion; it does not license a stronger conclusion than the proof, finite case, replay row, or comparator can carry.

Synthesis checkpoint 8. The preceding entries should be read together: promotion is earned by an alignment between claim wording, evidence class, and reopening rule. If any one of those three moves, the public sentence has to move with it.

\subsection{Physical constant replay: speed-of-light route}
The scientific reading is as follows. Evidence support: calibration replay of an official constant, not a new law of physics numeric replay QA with baseline, negative control, and falsifier repository path validation/numeric\_replay\_qa/OC133\_NUMERIC\_REPLAY\_QA\_TABLE.json This is replay QA and falsifier plumbing, not empirical theory promotion. The release consequence is local: this entry supports only the stated claim under its declared assumptions and can be reopened at the named boundary.

\subsection{Chemistry replay: water mass from a public source snapshot}
The evidence interpretation is deliberately bounded. Evidence support: NIST Chemistry WebBook molecular-weight field replay for water only numeric replay QA with baseline, negative control, and falsifier repository path validation/numeric\_replay\_qa/OC133\_NUMERIC\_REPLAY\_QA\_TABLE.json This is replay QA and falsifier plumbing, not empirical theory promotion. The important distinction is between the public statement, the artifact that carries it, and the condition that would make the statement too strong.

\subsection{Chemistry replay: water formula reconstruction}
The reviewer route starts from support, not from assertion. Evidence support: PubChem molecular-weight field replay for water only numeric replay QA with baseline, negative control, and falsifier repository path validation/numeric\_replay\_qa/OC133\_NUMERIC\_REPLAY\_QA\_TABLE.json This is replay QA and falsifier plumbing, not empirical theory promotion. A hostile reader should test the support class first, then ask whether the wording stays within that support class.

\subsection{Biology and geodata replay: counted public records}
The boundary reading prevents overclaim. Claim boundary: official GEO query-count replay; no organism-wide mechanism claim Evidence support: numeric replay QA with baseline, negative control, and falsifier repository path validation/numeric\_replay\_qa/OC133\_NUMERIC\_REPLAY\_QA\_TABLE.json This is replay QA and falsifier plumbing, not empirical theory promotion. This entry is included because it constrains promotion; it does not license a stronger conclusion than the proof, finite case, replay row, or comparator can carry.

Synthesis checkpoint 12. The preceding entries should be read together: promotion is earned by an alignment between claim wording, evidence class, and reopening rule. If any one of those three moves, the public sentence has to move with it.

\subsection{Systems replay: public macro-series route}
The scientific reading is as follows. Claim boundary: retrospective WDI GDP snapshot replay QA with descriptive comparator, no prediction or unrestricted comparative claim Evidence support: numeric replay QA with baseline, negative control, and falsifier repository path validation/numeric\_replay\_qa/OC133\_NUMERIC\_REPLAY\_QA\_TABLE.json This is replay QA and falsifier plumbing, not empirical theory promotion. The release consequence is local: this entry supports only the stated claim under its declared assumptions and can be reopened at the named boundary.

\subsection{Mathematical finite-model replay route}
The evidence interpretation is deliberately bounded. Evidence support: finite witness acceptance count for machine-checked current formal profile theorem inventory rows; never public promotion, empirical, or prediction support in current formal profile unless a future prospective protocol explicitly flips release promotion permission, prediction support permission, and empirical support permission after review numeric replay QA with baseline, negative control, and falsifier repository path validation/numeric\_replay\_qa/OC133\_NUMERIC\_REPLAY\_QA\_TABLE.json This is replay QA and falsifier plumbing, not empirical theory promotion. The important distinction is between the public statement, the artifact that carries it, and the condition that would make the statement too strong.

\subsection{Novelty boundary and residual-delta comparison}
The reviewer route starts from support, not from assertion. Claim boundary: No uniqueness, priority, absence, or invention claim is promoted by this row; systematic search remains future work. Evidence support: repository path comparators/OC\_1\_3\_3\_NOVELTY\_AND\_PRIORITY\_REGISTER.json Recorded content: Prior-art comparison is an illustrative bounded positioning note; uniqueness, priority, and absence are not promoted until a systematic search exists. ILLUSTRATIVE\_PRIOR\_ART\_POSITIONING\_ONLY\_NO\_UNIQUENESS\_PROMOTION A hostile reader should test the support class first, then ask whether the wording stays within that support class.
